% !PRINTED 405
\weberpage{405}
%<<CONT
\gesperrt{of such a kind that $u, u_1, \ldots, u_{\nu-1}$ can be expressed
rationally through $\vartheta$ alone.}

To prove it, we set

\weq{1}{u = \dfrac{\varphi(t)}{\psi(t)}, \quad u_1 = \dfrac{\varphi_1(t)}{\psi_1(t)}, \quad \cdots, \quad u_{\nu-1} = \dfrac{\varphi_{\nu-1}(t)}{\psi_{\nu-1}(t)},}

and understand by $\varphi(t)$, $\psi(t)$ integral functions of $t$
without common divisor; likewise by $\varphi_1(t)$,
$\psi_1(t)$, and so on. Moreover we may further assume that
none of the functions $u, u_1, \ldots, u_{\nu-1}$ is a constant.
Alongside the variable $t$ we introduce a second variable $\tau$ and
set:

\weq{2}{v = \dfrac{\varphi(\tau)}{\psi(\tau)}, \quad v_1 = \dfrac{\varphi_1(\tau)}{\psi_1(\tau)}, \quad \cdots, \quad v_{\nu-1} = \dfrac{\varphi_{\nu-1}(\tau)}{\psi_{\nu-1}(\tau)}.}

Now we form the integral functions of the two variables $t, \tau$,
of which none vanishes identically:

\weq{3}{\begin{aligned}
\varphi(t)\,\psi(\tau) - \varphi(\tau)\,\psi(t) &= \Phi(t,\tau) = -\Phi(\tau,t)\\
\varphi_1(t)\,\psi_1(\tau) - \varphi_1(\tau)\,\psi_1(t) &= \Phi_1(t,\tau) = -\Phi_1(\tau,t)\\
&\cdots\cdots\cdots\cdots\cdots\cdots\cdots\cdots\cdots\cdots\cdots\cdots\cdots\cdots\\
\varphi_{\nu-1}(t)\,\psi_{\nu-1}(\tau) - \varphi_{\nu-1}(\tau)\,\psi_{\nu-1}(t) &= \Phi_{\nu-1}(t,\tau) = -\Phi_{\nu-1}(\tau,t).
\end{aligned}}

If we regard them as functions of $t$, they all vanish
when $t = \tau$, and hence they have a greatest
common divisor, which with respect to $t$ is at least of the
first degree, and which shall be of degree $\mu$:

\weq{4}{R = R(t, \tau).}

This function $R$ is rational also with respect to $\tau$, and if
one arranges it so that $\tau$ occurs neither in the denominator nor in a
superfluous factor, then $R(t, \tau)$ can at most change its sign
under the interchange of $t$ and $\tau$.
For the functions (3), regarded also as functions of the two
variables $t, \tau$, are divisible by $R(t, \tau)$, and, since they
are alternating, by $R(\tau, t)$ (Bd.\,I, \S.\,51). Hence
$R(t, \tau)$ is divisible by $R(\tau, t)$, and conversely, and
consequently the two differ only by a constant factor,
which must be $= \pm 1$, because interchanging $t$ with
$\tau$ once more restores the original function.

It can further be proved that none of the functions (3) (for
indeterminate $\tau$), regarded as a function of $t$, contains a factor
%CONT>>
